% FIG ch13-01 pixels, physical size, effective resolution (Chinese)
% 栏宽按 1:2 绘制（栏宽 1 mm = 本页 0.5 mm），与英文图同一版式。
\begin{tikzpicture}[x=1mm, y=1mm, font=\sffamily\fontsize{8}{10}\selectfont, text=BookInk,
    >=Stealth]
  \tikzset{
    lbl/.style={inner sep=0pt, outer sep=0pt, anchor=south west},
    result/.style={inner sep=0pt, outer sep=0pt, anchor=north west},
    tick/.style={inner sep=0.4pt, outer sep=0pt, anchor=north},
    takeaway/.style={draw=BookTeal, line width=0.6pt, fill=white, rounded corners=1.2pt,
      inner sep=1.8mm, align=center, text width=122mm, font=\sffamily\fontsize{8}{10.5}\selectfont},
  }

  % 同一个位图文件
  \draw[BookRule, line width=0.6pt, fill=BookPale, rounded corners=1.2pt]
    (0.6, 43.0) rectangle (30.0, 76.8);
  \begin{scope}
    \clip (6.2, 63.2) rectangle (24.2, 73.325);
    \fill[white] (6.2, 63.2) rectangle (24.2, 73.325);
    \foreach \i in {0,...,15} {
      \foreach \j in {0,...,8} {
        \pgfmathparse{mod(\i+\j,2) ? 18 : 8}
        \fill[BookTeal!\pgfmathresult]
          (6.2+1.125*\i, 63.2+1.125*\j) rectangle ++(1.125, 1.125);
      }
    }
  \end{scope}
  \draw[BookRule, line width=0.4pt] (6.2, 63.2) rectangle (24.2, 73.325);
  \node[align=center, inner sep=0pt, anchor=center] at (15.3, 53.2)
    {{\fontsize{9}{11}\selectfont 文件}\\[2.5pt]1920 × 1080\,像素};

  % 两条宽度条的共同起点。85 mm -> 42.5 mm；180 mm -> 90 mm。
  \def\xO{36.2}
  \def\sS{0.5}

  \node[lbl] at (\xO, 73.4) {单栏 85\,mm};
  \fill[BookTeal] (\xO, 68.4) rectangle ({\xO+85*\sS}, 71.8);
  \node[result] at (\xO, 66.6) {1920 ÷ (85 ÷ 25.4) ≈ \textbf{574\,ppi}};

  \node[lbl] at (\xO, 56.6) {双栏 180\,mm};
  \fill[BookTeal] (\xO, 51.6) rectangle ({\xO+180*\sS}, 55.0);
  \node[result] at (\xO, 49.8) {1920 ÷ (180 ÷ 25.4) ≈ \textbf{271\,ppi}};

  \draw[BookInk, line width=0.55pt, ->] (30.0, 70.1) -- (\xO, 70.1);
  \draw[BookInk, line width=0.55pt, ->] (30.0, 53.3) -- (\xO, 53.3);

  % 与宽度条同一比例的标尺
  \draw[BookInk, line width=0.7pt] (\xO, 41.2) -- ({\xO+180*\sS}, 41.2);
  \foreach \mm/\lab in {0/0, 85/85, 180/180\,mm} {
    \draw[BookInk, line width=0.55pt] ({\xO+\mm*\sS}, 41.2) -- ({\xO+\mm*\sS}, 39.6);
    \node[tick] at ({\xO+\mm*\sS}, 39.3) {\lab};
  }
  \node[inner sep=0pt, outer sep=0pt, anchor=north west] at (0.6, 34.6)
    {同一比例尺（栏宽 1\,mm = 本页 0.5\,mm）};

  \node[takeaway, anchor=north west] at (0.6, 29.2) {%
    有效分辨率 = 像素数 ÷（宽度 mm ÷ 25.4）\\[2pt]
    改动文件里的 dpi 标签不改变像素数};
\end{tikzpicture}
